\subsection{Cross-Asset Replication and the Floor Effect}
\label{sec:results:cross}

Table~\ref{tab:cross} and Figure~\ref{fig:cross} repeat the analysis for QQQ and IWM. The fixed dollar thresholds of Section~\ref{sec:methods:criteria} can produce a floor effect for smaller underlyings, which we report; the scaled series of Section~\ref{sec:methods:arms} show whether the model's verdict still follows the criteria when the floor is lifted.
% src: report.yaml diagnostics.floor_effect (max_window_mean_abs_gex_bn) and diagnostics.floor_effect_note; detection_by_year.csv scope=all (QQQ, IWM, SPY); windows.csv.gz (clf_persistence_pct, clf_sign_flips, clf_avg_magnitude)
The floor effect is complete. Under the fixed thresholds the rule detects no window of QQQ (0 of 3{,}867, windows ending 2011 to 2026-10-09) and none of IWM (0 of 4{,}678, 2008 to 2026-10-09), in any year and on either series. The largest 30-day mean $|\mathrm{GEX}|$ is \$4.92~billion for QQQ and \$4.98~billion for IWM on \ident{new_full} (\$4.11~billion and \$4.37~billion without same-day expiries), against \$16.00~billion for SPY, so no window reaches the \$5~billion criterion. This is a consequence of the fixed thresholds, a design decision fixed before the runs, and not a modelling result. The magnitude criterion is the binding one: the persistence and sign-flip tests alone pass 1{,}434 of 3{,}867 windows of QQQ (37.1\%) and 3{,}357 of 4{,}678 of IWM (71.8\%), against 1{,}763 of 4{,}671 windows of SPY (37.7\%). gpt-6-luna labels no window of QQQ or IWM as persistent either (0 of 3{,}867 and 0 of 4{,}678 on \ident{new_full}; 0 of 1{,}699 and 0 of 1{,}700 on \ident{new_ex0dte}), so the unscaled arms show that the model also rejects every window; they cannot test its discrimination.

% src: report.yaml diagnostics.scaling (factor 4.224724 and 11.43094 shown as 4.22 and 11.43, n_days_in_factor 1699 and 1700); series parquet new_full_scaled_QQQ / _IWM series_metadata (ratio_min, ratio_max, ratio_p25, ratio_p75); stats.yaml contrasts series_contrast family scale and period_shift for new_full_scaled
The scaled variant multiplies the \ident{new_full} values by one factor per symbol: 4.22 for QQQ (median of 1{,}699 daily ratios) and 11.43 for IWM (1{,}700 days). It is a disclosed protocol change (Section~\ref{sec:methods:arms}). With it the rule detects 350 of 1{,}670 windows of QQQ (21.0\%) and 901 of 1{,}671 windows of IWM (53.9\%), and gpt-6-luna 360 (21.6\%) and 904 (54.1\%). The model follows the rule on the scaled inputs: the verdicts agree in 99.4\% of the QQQ windows (kappa 0.982) and 99.7\% of the IWM windows (kappa 0.994). The paired difference to the unscaled arms is 21.0 percentage points for QQQ (quarter-block bootstrap interval [12.6, 31.2], $p=5.0\times10^{-4}$; exact McNemar test on the non-overlapping subsample $p=4.9\times10^{-4}$) and 53.9 for IWM ([40.7, 67.4], $p=5.0\times10^{-4}$; $p=4.7\times10^{-10}$); for the model the differences are 21.6 and 54.1 percentage points.

The scaled rates depend on the factor and are not a measurement of QQQ or IWM regimes. The ratio of the open-interest notional of SPY to that of the symbol is not constant: for QQQ the daily ratio runs from 0.78 to 38.32 (quartiles 3.30 and 5.73), and for IWM from 1.51 to 110.57 (quartiles 7.72 and 16.46). A single factor puts the exposure on the scale of SPY on a typical day, no more, and another factor would give other rates. The scaled arms therefore show only that the model follows the criteria once the floor is lifted. They show no step after 2023: the rate of the rule on the scaled IWM series is 65.2\% in 2020 to 2023 and 35.7\% in 2024 and 2025 (shift $-$29.5 percentage points, $p=0.061$ by bootstrap, $p=0.137$ by Fisher exact test on the subsample), and on the scaled QQQ series 25.9\% and 13.4\% ($-$12.5 percentage points, $p=0.170$, $p=0.292$).

\begin{table}[H]
\centering
\caption{Cross-asset replication: windows, typical magnitude and detected windows under the fixed thresholds, by symbol and series. The scaled series are the disclosed protocol change of Section~\ref{sec:methods:arms}. The magnitude column is the median over windows of the 30-day mean $|\mathrm{GEX}|$ of the series as sent to the model; the scaled series are SPY-scale values, so their magnitudes are not comparable with the unscaled rows of the same symbol. The ex-0DTE arms of QQQ and IWM start in 2020, so gpt-6-luna answered fewer windows than the classifier labelled.}
\label{tab:cross}
\footnotesize
\begin{tabular}{@{}llrrrr@{}}
\toprule
Symbol & Series & Windows & \shortstack[r]{Median mean\\ $|\mathrm{GEX}|$ (\$B)} & \shortstack[r]{Detected\\ (classifier)} & \shortstack[r]{Detected\\ (gpt-6-luna)} \\
\midrule
% src: detection_by_year.csv scope=all symbol=SPY (classifier series_key=new_full: n, n_detect; llm arm=full_spy: n_detect); windows.csv.gz series_key=new_full symbol=SPY (median of clf_avg_magnitude)
SPY & \ident{new_full} & 4{,}671 & 3.84 & 606 of 4{,}671 (13.0\%) & 613 of 4{,}671 (13.1\%) \\
% src: detection_by_year.csv scope=all symbol=QQQ (classifier series_key=new_full: n, n_detect; llm arm=full_qqq: n_detect); windows.csv.gz series_key=new_full symbol=QQQ (median of clf_avg_magnitude)
QQQ & \ident{new_full} & 3{,}867 & 0.74 & 0 of 3{,}867 (0.0\%) & 0 of 3{,}867 (0.0\%) \\
% src: detection_by_year.csv scope=all symbol=IWM (classifier series_key=new_full: n, n_detect; llm arm=full_iwm: n_detect); windows.csv.gz series_key=new_full symbol=IWM (median of clf_avg_magnitude)
IWM & \ident{new_full} & 4{,}678 & 0.66 & 0 of 4{,}678 (0.0\%) & 0 of 4{,}678 (0.0\%) \\
% src: detection_by_year.csv scope=all symbol=SPY (classifier series_key=new_ex0dte: n, n_detect; llm arm=ex0dte_spy: n_detect); windows.csv.gz series_key=new_ex0dte symbol=SPY (median of clf_avg_magnitude)
SPY & \ident{new_ex0dte} & 4{,}671 & 3.04 & 491 of 4{,}671 (10.5\%) & 494 of 4{,}671 (10.6\%) \\
% src: detection_by_year.csv scope=all symbol=QQQ (classifier series_key=new_ex0dte: n, n_detect; llm arm=ex0dte_qqq: n_detect); windows.csv.gz series_key=new_ex0dte symbol=QQQ (median of clf_avg_magnitude)
QQQ & \ident{new_ex0dte} & 3{,}867 & 0.64 & 0 of 3{,}867 (0.0\%) & 0 of 1{,}699 (0.0\%) \\
% src: detection_by_year.csv scope=all symbol=IWM (classifier series_key=new_ex0dte: n, n_detect; llm arm=ex0dte_iwm: n_detect); windows.csv.gz series_key=new_ex0dte symbol=IWM (median of clf_avg_magnitude)
IWM & \ident{new_ex0dte} & 4{,}678 & 0.62 & 0 of 4{,}678 (0.0\%) & 0 of 1{,}700 (0.0\%) \\
% src: detection_by_year.csv scope=all symbol=QQQ (classifier series_key=new_full_scaled: n, n_detect; llm arm=full_qqq_scaled: n_detect); windows.csv.gz series_key=new_full_scaled symbol=QQQ (median of clf_avg_magnitude)
QQQ & \ident{new_full_scaled} & 1{,}670 & 9.13 & 350 of 1{,}670 (21.0\%) & 360 of 1{,}670 (21.6\%) \\
% src: detection_by_year.csv scope=all symbol=IWM (classifier series_key=new_full_scaled: n, n_detect; llm arm=full_iwm_scaled: n_detect); windows.csv.gz series_key=new_full_scaled symbol=IWM (median of clf_avg_magnitude)
IWM & \ident{new_full_scaled} & 1{,}671 & 12.00 & 901 of 1{,}671 (53.9\%) & 904 of 1{,}671 (54.1\%) \\
\bottomrule
\end{tabular}
\end{table}

\begin{figure}[H]
\centering
\resultfigure{\figCross}
\caption{Cross-asset replication, SPY, QQQ and IWM (QQQ from 2011). Top row: detection rate by year of the classifier (lines) and gpt-6-luna (markers) on \ident{new_full} and \ident{new_ex0dte}, and the share of windows that pass the persistence and sign-flip tests alone (dotted). Bottom row: yearly median and largest 30-day mean $|\mathrm{GEX}|$ against the \$5 billion criterion. \textbf{Read this figure as:} for QQQ and IWM the detection lines lie at zero in every year while the dotted lines show that many windows pass the persistence and flip tests; the bottom row shows why, since the largest yearly window mean stays below the \$5 billion line for QQQ (at most \$4.92~billion) and for IWM (at most \$4.98~billion). For SPY the largest window mean reaches \$16.00~billion.}
\label{fig:cross}
\end{figure}
